% Nouveau EGA V (cote 25): archive cote 25, site page 88 (PDF page 89).
% typescript. Draft compared: none.
% AI-transcribed by Claude Opus 5.5 (claude-opus-5-5, adaptive thinking, effort high), Message Batches API, batch ega25b (msgbatch_01A7n5YuXubUPNK1zsepi8RG), 2026-11-13; not yet checked by eye.
\scanpage{25/88}
\mspage{\hand{82}}

Grass$^{m}$(P), la dimension de $\underline{X}^{(m)}_{\hand{\xi}}$ est au moins égale à dim X-m , il en est
% NOTE: the right edge of the scan cuts the ends of the first two lines ("est", "l'inéga")
en particulier ainsi lorsque $\hand{\xi}$ provient d'un point de Q(C$_{\hand{y}}$). Pour l'inéga
lité en sens inverse, notons que (nous plaçant sur le corps k' = k($\hand{\xi}$))
comme $C_{\hand{y}\,k'} \subset L_{\hand{\xi}}$ est \xout{de} un hyperplan de $L_{\hand{\xi}}$ , si la dimension de
$\underline{X}^{(m)}_{\hand{\xi}} = X\times_{P}L_{\hand{\xi}}$ était \over{>}{$\geqslant$} dim X - m + 1 , alors celle de $\underline{X}^{(m+1)}_{\hand{y}}$ =
$X\times_{P}C_{\hand{y}}$ serait \over{>}{$\geqslant$} dim X - m \xout{+}, (puisque le changement de base k($\hand{y}$)$\to$k'
trans\xout{g}forme ce dernier préschéma en $(X\times_{P}L_{\hand{\xi}})\times_{L_{\hand{\xi}}}(C_{\hand{y}\,k'})$ ). Or on a au
contraire \xout{dimxXx$_{P}$C} dim $\underline{X}^{(m+1)}$ = dim X - m -1 \xout{(avecxlaxconventionxque}
\deleted{(du moins en supposant dim f(X) \hand{$>$} m )}, en vertu du n°2 (revu au n°10).
Le \inserted{cas b)} \deleted{cas où dim f(X)=m est justiciable également du cas b), qui} se traite
de façon analogue: si on avait dim $X\times_{P}L$ \hand{$\geqslant$} d\hand{+1}, \xout{alorsxxxlimagexde} ou
ce qui revient au même \xout{f(X)xxf$_{k(\ )}$xxxx} $f_{k'}(X_{k'}) \cap L_{\hand{\xi}}$ de dimension \hand{$\geqslant$}\over{0}{1},
alors on aurait par le même argument que plus \xout{haut} haut que
$\underline{X}^{(m+1)}_{\hand{y}} \neq \emptyset$ , contrairement à ce qu'on avait signalé avant 14.1.

\uline{Corollaire} 14.2. Supposons X de dimension m et f:X$\to$P fini\inserted{(resp quasi-fini)}, alors
le morphisme $p_{C_{\hand{y}}} : X_{k(\hand{y})} \longrightarrow Q(C_{\hand{y}})$ est fini surjectif\hand{(}resp. quasi-fini
dominant).

En effet, ce morphisme est \inserted{quasi-}fini, et comme dim $X_{\hand{k(y)}}$ = dim Q (C$_{\hand{y}}$), il est
dominant; \xout{enfinxdoncxsurjectif} si f est fini, $p_{C_{\hand{y}}}$ est de plus fini,
\xout{doncxfex} car propre, donc surjectif étant dominant.
\uline{Corollaire} 14.3. Sous les conditions de 14.1. a), si X est Cohen-Macau-%
lay, alors le morphisme $p_{C} : \widetilde{X}(C_{\hand{y}}) \longrightarrow Q(C_{\hand{y}})$ est un morphisme de
Cohen-Macaulay, et à fortiori \uline{plat}.

Pour la démonstration, cf remarque en haut de la page 21, avant 5.6.\inserted{qui donne un résultat plus fort, (à inclure dans 14.3), savoir que les $\unsure{\varphi^{-1}_{\xi}}$ pour $\xi\in Q(C_{y})$ sont \unsure{$F_{\xi}$}-réguliers.}
Ce \over{c}{C}orollaire devrait être modulé, mais pour simplifier on pourra
supposer f quasi-fini: si F est un Module de Cohen-Macaulay sur X, et si
pour to\uline{u}te composante irréductible Z de Supp F, on a dim Z \over{>}{$\geqslant$} m, alors
$\widetilde{F}(C_{\hand{y}})$ est CM \hand{(}et à fortiori plat\hand{)} relativement à Q(C$_{\hand{y}}$).

On notera qu'on ne peut pas remplacer,\xout{xxxx} pour obtenir la

\margin{88}
